\Qhead{4.2}{Оптические резонаторы}{Оптические резонаторы. Резонатор Фабри-Перо, конфокальный и концентрический резонаторы. Неустойчивый резонатор. Продольные и поперечные типы колебаний. Спектр частот и расходимость излучении. Добротность.}
\Subfile{4.2.1}{sections/04_generation/q02/s01}

\recap{Резонатор задаёт пространственное самовоспроизведение и фазовый спектр. Продольные частоты определяет время обхода, поперечные — фаза Гуи; потери задают время жизни и ширину.}
\board{Фабри–Перо; конфокальная и концентрическая геометрии; $Q=\omega\tau_p$, FSR и финесс.}
\oral{Почему ширина лазерной линии может быть существенно меньше ширины пассивного резонанса?}
\bait{Можно оставить фазу Гуи: она различает частоты поперечных мод с разным числом узлов.}
\sources{Митрофанов~\cite{mitrofanov}, лекция~13, PDF-с.~60--65; Kärtner~\cite{kartnerfull}, глава~2. Не путать ширину резонанса и шумовую линию генерации следующего вопроса.}
